% ===========================================================================
\chapter{Where to go next}\label{ch:next}

The agent is small: about 2,300 lines in ten files, comments and all. The coding
agents people use every day do far more, but most of what they add is
built on what you have here: the loop, the tools, and the rules around
them. Here are some directions, roughly from easiest to hardest.

\section{Small additions}

\textbf{Project instructions.} Many projects keep notes for agents in a
file such as \code{AGENTS.md}: how to run the tests, the house style,
what not to touch. Read it at the start, and add it to the system prompt.

\textbf{Cost.} \code{Usage} counts tokens. Multiply by the model's prices
(fresh input, cached input and output are priced differently) and print
what each task cost.

\textbf{Undo.} Before every edit or write, keep the file as it was. An
\code{undo} command (or a tool) puts it back. Simpler still, in a
project under git: commit before each task, so that \code{git diff} shows
what the agent did and \code{git checkout .} takes it all back.

\textbf{More tools.} A tool to delete or rename a file; one to look at a
web page; one that runs just the tests for one file. Each is a class
inheriting \code{Tool}, plus checks. Write its description as carefully
as its code.

\section{Doing things at once}

\textbf{Several tool calls in parallel.} A model often asks for three files
at once. The agent runs them one after another; with \code{spawn}, it
could read them at the same time. Only for tools that change nothing,
though: two edits to one file must not race.

\textbf{Streaming.} The agent waits for each reply to be complete before
showing any of it. The API can instead \emph{stream} a reply, sending it as
it is written, as a series of events; the agent can then show the model's
text as it appears. With streaming, \code{max_tokens} can be much larger,
too, since a long reply no longer risks an HTTP timeout. It is mostly HTTP
work: reading the response line by line, and assembling the blocks from
the events.

\section{Bigger ideas}

\textbf{Sub-agents.} When a task needs a lot of reading (``find everything
that uses the old API''), all that reading fills the conversation. Give
the agent a tool that starts \emph{another} agent, with its own fresh
conversation and only the reading tools, to do the search and report back
in a paragraph. The main conversation only ever sees the paragraph.

\textbf{A plan.} For a long task, a tool with which the model keeps a
checklist (``1. find the bug; 2. fix it; 3. add a check'') and ticks items
off. It helps the model keep track, and it shows you where it is.

\textbf{Server-side compaction.} Chapter~\ref{ch:long}'s summing up, done by
the API instead, in a form designed for it.

\textbf{A stronger sandbox.} The workspace rule covers the file tools; a
command can still do anything you can. Run the commands in a container
with only the workspace mounted, and without the network, and the
permission questions can be fewer.

\section{Testing an agent}

Our checks test the agent's code. Whether it does its \emph{job} well is a
different question, and a statistical one: a model can solve a task four
times out of five. To answer it, collect tasks with a way to check the
result (here: ``do the checks pass afterwards?''), run each several
times, and count. Then every change to a tool description or to the system
prompt can be measured, not guessed at. Recorded sessions help here too:
when a run goes wrong, its recording shows exactly what the model saw.
